% =====================================================================
%  tr01a —— §3 讲义（原文 pp.8--9）
%  归属：tr-01。含 §3 标题、"General background"、Teichmüller 空间与映射类群
%  的开场、Fenchel--Nielsen 坐标；编号陈述：定理 1（p.9）；图 2（p.8）、图 3（p.9）。
%  原文本块内无编号公式（编号公式只有 p.18 的 (1)(2)(3) 与 p.21 的 (4)）。
% =====================================================================

% [原文 p.8]
\section{讲义}

下面是对 Mirzakhani 关于 Witten--Kontsevich 定理之证明的阐述，其中包含关于
Teichm\"uller 理论、模空间理论以及辛约化的直接背景材料。这些讲义是按研究生教材的
方式呈现的——整体阐述沿着概念的发展展开，而不考虑材料的历史发展脉络。目标是给出
材料的整体处理，并使读者获得整体性的理解。在若干地方，为了对材料作更简化的处理，
把好几位作者的途径结合了起来。文中只收录了与主题直接相关的文献。读者应当去查阅
文献，以了解历史发展、完整的参考文献以及材料的各种推论。

\begin{enumerate}[leftmargin=1.8em,itemsep=2pt,topsep=2pt]
\item \textbf{背景与概览。}
\item \textbf{McShane--Mirzakhani 恒等式。}
\item \textbf{余体积公式与递推。}
\item \textbf{辛约化、$S^1$ 主丛与正规形式。}
\item \textbf{相交数的模式与 Witten--Kontsevich。}
\end{enumerate}

\para{第 1 讲：背景与概览。}

\para{一般背景。}

由单值化（Uniformization），对负 Euler 示性数的曲面而言，一个共形结构等价于一个
完备的双曲结构。我们考虑具有双曲度量、有限拓扑型的 Riemann 曲面 $R$：若边界非空，
则允许带穿刺点与测地边界。固定一个拓扑参考曲面 $F$，并考虑一个标记（marking），
即一个保向同胚 $f:F\to R$，两个这样的同胚在相差一个同伦的意义下视为等同。若边界非空，
同伦是保持边界（作为集合）意义下的。
记 $\T$ 为 $R$ 的 Teichm\"uller 空间——即等价类对 $\{(f,R)\}$ 所构成的空间，其中
两个对等价，是指存在一个同伦意义下的映射三角形，其第三条边是 Riemann 曲面之间的
共形映射（双曲等距）。

% [Tutor Note: p.8 的"标记（marking）"与 Teichmüller 点 $\{(f,R)\}$ 的等价关系要
%  读准：$f$ 只到同伦为止（边界非空时到 rel boundary 的同伦为止），所以 $\T$ 忘掉
%  标记后得到的正是模空间。维数锚点：$\T_{g,n}$ 的复维数是 $3g-3+n$，实维数
%  $6g-6+2n$，与定理 1 中 FN 参数恰有 $3g-3+n$ 个（每个贡献 $\ell_j$ 与 $\tau_j$
%  两个实参数）完全吻合；反之若有人把 $\ell$ 或 $\tau$ 之一漏掉，维数立刻对不上。]

% [原文 p.8] 图 2
\begin{figure}[htbp]
\centering
\includegraphics[width=0.75\textwidth]{fig02.png}
\figtitle{2}{一个亏格 3、带 2 条测地边界的曲面。}
\end{figure}

我们考虑下列 Teichm\"uller 空间 $\T$：$\Tg$——对紧的亏格 $g$ 曲面；
$\Tg(L)$——对亏格 $g$、带标号的测地边界且边界长度预先给定的曲面；
$\Tgn$——对亏格 $g$、带 $n$ 个标号穿刺点的曲面。

% [原文 p.9]
在 $\Tg(L)$ 的情形，曲面之间的同伦是保持边界（作为集合）意义下的。$\Tg$ 与 $\Tgn$ 是复流形，
而 $\Tg(L)$ 是实解析流形。

$F$ 上一个非平凡、非穿刺点边缘（non puncture peripheral）的自由同伦类
$\alpha$，在曲面 $R$ 上对 $f(\alpha)$ 有唯一的测地代表——测地长度 $\ell_\alpha(R)$
于是给出 Teichm\"uller 空间上的一个自然函数。测地长度函数的集合提供局部坐标，并给出
到 Euclid 空间的整体浸入。对简单曲线而言，测地长度的微分处处不为零。在 $\T$ 的每
一点，简单曲线的测地长度微分在余切丛中稠密。

% [Tutor Note: p.9 这段把"测地长度函数给出局部坐标"与"整体浸入"分开说，是两件
%  不同的事：局部是 $3g-3+n$ 个（复）坐标函数的坐标卡，整体只是到
%  $\R^{6g-6+2n}$ 的浸入而非嵌入（$\T$ 非紧，且长度函数在趋于退化曲面时趋于 $0$）。
%  "简单曲线的测地长度微分在余切丛中稠密"这一条是 Wolpert 文献里反复用到的关键
%  手段：它保证两个 $2$-形式只要在所有 $\dd\ell_\alpha$（$\alpha$ 取遍简单闭曲线）
%  上取值相同，就处处相同；后文用 $\dd\ell$ 检验辛形式（如定理 2）时靠的正是这条。]

一个曲面可以沿一条简单闭测地线剪开——边界成为等距圆周。由于简单测地线的一个邻域
具有 $S^1$ 对称性，可以把边界作相对旋转后重新拼合，得到新的双曲结构。该形变就是
Fenchel--Nielsen（FN）扭转。

% [原文 p.9] 图 3
\begin{figure}[htbp]
\centering
\includegraphics[width=0.6\textwidth]{fig03.png}
\figtitle{3}{一个正的 Fenchel--Nielsen 扭转形变。}
\end{figure}

对初始相邻点作单位速度双曲位移所得的无穷小形变，是 $\T$ 上的 Fenchel--Nielsen
无穷小扭转向量场 $t_\alpha$。（正扭转对应于穿过该测地线时向右位移。）

同长度的双曲曲面测地边界可以拼合起来构成新曲面。给定边界参考点之后，相对旋转以弧长
度量。一条双曲裤子（pair of pants）是带三条测地边界的亏格零曲面。对裤子而言，边界
参考点由边界之间唯一的正交连接测地线给出。在粘合处，公共边界长度 $\ell$ 与相对旋转，
即 FN 扭转参数 $\tau$，都是不受限制的参数（$\tau$ 由初始构形出发作延拓来定义）。
长度 $\ell$ 在 $\R_{>0}$ 中变化，扭转 $\tau$ 在 $\R$ 中变化。每个有限拓扑型的双曲
曲面都可以由裤子拼装而成。

\begin{thm}{1} Fenchel--Nielsen 坐标。固定一个裤子分解的拓扑型与一个初始构形，则 FN
参数 $\prod_{j=1}^{3g-3+n}(\ell_j,\tau_j)$ 定义一个从 $\T$ 到
$\prod_{j=1}^{3g-3+n}\R_{>0}\times\R$ 的实解析等价。\end{thm}

% [Tutor Note: p.9 的无穷小扭转向量场 $t_\alpha$ 与 p.10 的扭转–长度对偶
%  $\omega(\ ,t_\alpha)=\dd\ell_\alpha$ 是同一件事的两面。在 FN 坐标下取
%  $t_{\alpha_j}=\partial/\partial\tau_j$，代入 p.10 定理 2 的
%  $\omega=\sum_j\dd\ell_j\wedge\dd\tau_j$ 恰好得
%  $\omega(\,\cdot\,,t_{\alpha_j})=\dd\ell_j$，这是可以当场手算的最短核对。
%  注意 $t_\alpha$ 按"单位速度双曲位移"归一化；换成别的归一化，对偶式会差一个
%  常数因子，而 $\dd\ell\wedge\dd\tau$ 的系数也会随之改变。]
